\section{Round-five referee closure: projective finite-source commutation and global closed duality}
\label{sec:r5-d1}

The global rate is not reconstructed from one local analytic pressure.
Instead, the model supplies a compatible family of finite-source pressures
on increasing balls, and the global convex dual is the closed supremum of
their local conjugates.

\subsection{The projective pressure family}

Let \(X_1\subset X_2\subset\cdots\) be a countable family of finite-dimensional
or separable Hilbert source spaces whose union separates the weighted path
state.  For \(R<\infty\), let \(B_R\subset X_R\) be the real ball on which the
model pressure \(Q_R\) is obtained by finite-source continuation.  On
overlaps,
\[
 Q_{R'}|_{B_R}=Q_R\qquad(R'\ge R).
\]
For a path state \(x\), put
\[
 I_R(x)=\sup_{\Theta\in B_R}
 \{\langle\Theta,x\rangle-Q_R(\Theta)+Q_R(0)\},
 \qquad
 I(x)=\operatorname{cl}\sup_R I_R(x).
\]

\begin{theorem}[Global rate from compatible finite-source balls]
\label{thm:r5-d1-global}
Assume the B2 full LDP and exponential compact containment in the weighted
state space.  Then its good rate equals \(I\).  Equivalently,
\[
 I(x)=
 \sup_{\Theta\in\bigcup_RB_R}
 \{\langle\Theta,x\rangle-Q(\Theta)+Q(0)\}
\]
followed by lower-semicontinuous closure.  No value of \(Q\) outside the
proved union of finite-source balls is used.
\end{theorem}

\begin{proof}
The Laplace upper bound gives \(I\ge\sup_RI_R\).  Conversely, B2's smooth
positive recovery approximates every finite-rate point by points exposed by
some bounded finite source; that source lies in one \(B_R\).  The Laplace
lower bound at the exposed approximants gives equality there, and action
convergence plus lower semicontinuity gives equality at the original point.
Goodness identifies the closed envelope and rules out an additional
noncompact branch.
\end{proof}

\subsection{Derivative and conditioning commutation}

For a finite-volume observable \({\mathcal X}_\varepsilon\), define
\[
 Q_{\varepsilon,R}(\Theta)
 =\frac1{\mu_\varepsilon}
 \log\mathbb E
 e^{\mu_\varepsilon\langle\Theta,{\mathcal X}_\varepsilon\rangle}.
\]

\begin{proposition}[Exact normalization and derivative limit]
\label{prop:r5-d1-normalization}
For \(F,G\in X_R\),
\[
 DQ_{\varepsilon,R}(\Theta)F
 =\mathbb E_{\Theta,\varepsilon}\langle F,{\mathcal X}_\varepsilon\rangle,
\]
\[
 D^2Q_{\varepsilon,R}(\Theta)(F,G)
 =\mu_\varepsilon
 \operatorname{Cov}_{\Theta,\varepsilon}
 (\langle F,{\mathcal X}_\varepsilon\rangle,
  \langle G,{\mathcal X}_\varepsilon\rangle).
\]
Local uniform holomorphic convergence on a larger complex ball implies local
uniform convergence of every fixed Fr\'echet derivative.
\end{proposition}

\begin{proof}
The finite identities follow by differentiation.  Cauchy's formula on the
larger complex ball gives uniform convergence of derivatives on each smaller
ball.  The statement is projective in \(R\), so compatible derivatives agree
on overlaps.
\end{proof}

Let \(C:X^*\to\mathbb R^k\) be a finite constraint with positive covariance.
At finite volume choose the exact saddle
\[
 C\,DQ_{\varepsilon,R}(\Theta+C^*\lambda_{\varepsilon,\Theta})=a.
\]

\begin{theorem}[Finite-saddle conditioning theorem]
\label{thm:r5-d1-conditioning}
The normalized conditioned pressure is
\[
 Q_{\varepsilon,R}^a(\Theta)
 =
 Q_{\varepsilon,R}(\Theta+C^*\lambda_{\varepsilon,\Theta})
 -\lambda_{\varepsilon,\Theta}\cdot a
 -Q_{\varepsilon,R}(C^*\lambda_{\varepsilon,0})
 +\lambda_{\varepsilon,0}\cdot a.
\]
It converges with all fixed derivatives, and
\[
 D^2Q_R^a
 =
 Q_{\Theta\Theta}
 -Q_{\Theta\lambda}Q_{\lambda\lambda}^{-1}
  Q_{\lambda\Theta}.
\]
The associated Gaussian covariance is the same Schur complement.
\end{theorem}

\begin{proof}
B1 supplies existence, uniqueness, and local uniform convergence of the exact
finite saddle.  The shell coefficient is subexponential and cancels between
numerator and denominator.  Differentiating the finite saddle equation gives
the envelope Hessian; Proposition~\ref{prop:r5-d1-normalization} identifies
it with the covariance of the scaled conditioned field.
\end{proof}

\subsection{Closed conditioning/contraction interchange}

Let \(T:X^*\to E\) be a typed continuous map or an exponentially good
approximation with continuous adjoint on sources.

\begin{theorem}[Closed affine interchange]
\label{thm:r5-d1-interchange}
The joint constrained contraction has rate
\[
 J_a(y)
 =
 \inf\{I(x):Tx=y,\ Cx=a\}.
\]
Its source representation is
\[
 J_a
 =
 \operatorname{cl}
 \left(Q^a\circ T^*\right)^*.
\]
If the affine image of every rate sublevel is closed, the closure symbol may
be omitted.  First and second pressure derivatives commute with \(T\) on
every finite source ball.
\end{theorem}

\begin{proof}
Apply the contraction theorem to the joint map \((T,C)\).  Fenchel duality
gives the conjugate of the constrained pressure, but projection of a closed
convex set need not be closed; taking the lower-semicontinuous closure is
therefore necessary in general.  Closedness of affine images of compact rate
sublevels removes it under the stated condition.  The finite-volume pressure
identity
\(Q_{\varepsilon,T}(\eta)=Q_\varepsilon(T^*\eta)\) is exact, and the chain
rule gives derivative commutation.
\end{proof}

\subsection{Finite-volume-centered Gaussian likelihoods}

Let
\[
 m_{\varepsilon,\Theta}=DQ_{\varepsilon,R}(\Theta),
 \qquad
 Z_\varepsilon
 =\sqrt{\mu_\varepsilon}
 ({\mathcal X}_\varepsilon-m_{\varepsilon,\Theta}).
\]

\begin{theorem}[Uniform local likelihood expansion]
\label{thm:r5-d1-lan}
For \(h\) in a compact subset of a finite-dimensional source subspace,
\[
 L_\varepsilon(h)
 =
 \exp\left\{
 \langle h,Z_\varepsilon\rangle
 -\mu_\varepsilon
 \left[
 Q_{\varepsilon,R}(\Theta+h/\sqrt{\mu_\varepsilon})
 -Q_{\varepsilon,R}(\Theta)
 -\frac1{\sqrt{\mu_\varepsilon}}
  DQ_{\varepsilon,R}(\Theta)h
 \right]\right\}
\]
converges jointly with \(Z_\varepsilon\) to
\[
 \exp\{\langle h,Z\rangle-\tfrac12\Sigma_\Theta(h,h)\}.
\]
The convergence is uniform on compact \(h\)-sets and the likelihoods are
uniformly integrable.  Under \(L_\varepsilon(h)\), the limit is shifted by
\(\Sigma_\Theta h\).
\end{theorem}

\begin{proof}
The source displacement lies in a fixed larger complex ball for all small
\(\varepsilon\).  Cauchy bounds give a uniform cubic remainder
\(O(|h|^3/\sqrt{\mu_\varepsilon})\).  The quadratic term converges by
Proposition~\ref{prop:r5-d1-normalization}; the characteristic-function
expansion gives the Gaussian limit.  A still larger real ball yields a
uniform \(L^{1+\delta}\) bound and hence uniform integrability.  Multiplying
the joint characteristic function by the limit likelihood gives the
Cameron--Martin shift.
\end{proof}

\subsection{Dynamic triangle}

Let \(I_{s,t}\) be the B2 action and \(S_{s,t}\) the B4 action semigroup.

\begin{theorem}[Rate--pressure--semigroup commutation]
\label{thm:r5-d1-triangle}
On every finite source ball and finite horizon,
\[
 Q_{s,t}\xleftrightarrow{\ \mathrm{Fenchel}\ }I_{s,t},
 \qquad
 (S_{s,t}\Phi)(f)
 =\sup_{\rho_s=f}\{\Phi(\rho_t)-I_{s,t}(\rho)\}.
\]
Differentiation gives the mean and covariance, the quadratic tangent of
\(I_{s,t}\) is the inverse-covariance form on the covariance range, and
dynamic additivity of \(I\) is equivalent to the nonlinear tower of \(S\).
These identities are compatible as \(R\) increases.
\end{theorem}

\begin{proof}
B2 gives pressure/action duality on the union of finite source balls and B4
constructs the variational semigroup.  Second-order expansion and
Theorem~\ref{thm:r5-d1-lan} identify the local quadratic dual.  Concatenating
balanced paths adds the action, while splitting a path gives the reverse
dynamic-programming inequality.  Overlap uniqueness of the pressure makes
the entire diagram projectively compatible.
\end{proof}

\subsection{Independent scalar calibration}

Let a nonconstant mechanical work \(G\) have differentiable contracted rate
\(I_G\) at target \(a\), and set \(\theta=I_G'(a)\).

\begin{proposition}[Canonical-cocycle calibration]
\label{prop:r5-d1-calibration}
The information projection onto \(\mathbb E G=a\) has density proportional to
\(e^{\theta G}\).  A cash-additive, strictly relevant, time-consistent
certainty equivalent generated by one \(C^2\) utility is entropic with a
constant coefficient \(\vartheta\).  Requiring its likelihood cocycle to be
the same mechanical information-projection cocycle forces
\(\vartheta=\theta\); without that requirement, \(\vartheta\) is independent.
\end{proposition}

\begin{proof}
Strict convex duality gives the unique scalar Lagrange multiplier
\(\theta=I_G'(a)\).  Translation invariance of a certainty equivalent gives
constant absolute risk aversion and hence affine or exponential utility; the
tower makes the exponential coefficient time independent.  Equality of the
two normalized densities on nonconstant \(G\) implies equality of their
exponents.
\end{proof}

\begin{theorem}[Round-five D1 closure]
\label{thm:r5-d1-main}
The deterministic theta diagram is a projective theorem on the proved union
of finite source balls: global rates use closed convex duality, conditioning
uses exact finite saddles, contraction includes the necessary closure,
Gaussian likelihoods use finite-volume centering, and the dynamic and
calibration diagrams commute on their typed domains.
\end{theorem}

\begin{proof}
Combine Theorems~\ref{thm:r5-d1-global},
\ref{thm:r5-d1-conditioning}, \ref{thm:r5-d1-interchange},
\ref{thm:r5-d1-lan}, and \ref{thm:r5-d1-triangle}, together with
Proposition~\ref{prop:r5-d1-calibration}.
\end{proof}
